% \iffalse meta-comment
%
% Copyright (C) 2017- by Yanshuo Chu <yanshuoc@gmail.com>
%
% This file may be distributed and/or modified under the
% conditions of the LaTeX Project Public License, either version 1.3a
% of this license or (at your option) any later version.
% The latest version of this license is in:
%
% http://www.latex-project.org/lppl.txt
%
% and version 1.3a or later is part of all distributions of LaTeX
% version 2004/10/01 or later.
%
% \fi
%
% \section{元信息字段的回落值}
%
%    \begin{macrocode}
%<*hithesis-cfg>
%    \end{macrocode}
%
%
% \changes{v3.2a}{2026/08/29}{元信息字段不填时的回落值照官方范例的占位形给：
%   姓名导师是方框、学号十个方框、题目取范例的题目。原先一律空着}
% \textbf{不填就印范例的占位形。} 官方范例本身就是一份「空白表」：姓名印
% \texttt{□□□}、导师印 \texttt{□□□教授}、学号印十个方框、题目是那条气体
% 静压轴承。回落值照抄它，什么都不填编出来的就是学校那份范例的样子——哪一格
% 忘了填，纸面上是一眼可见的方框，不是一段安静的空白。iota-hit 的
% \texttt{info-defaults} 同一张表。日期不在此列：那一项回落到编译当天（东八区、
% 年月），比范例里冻死的「2022 年 6 月」有用。
%    \begin{macrocode}
\ExplSyntaxOff
\hit@presetup{ info = {
  title-zh      = {局部多孔质气体静压轴承关键技术的研究},
  author-zh     = {□□□},
  supervisor-zh = {□□□教授},
  affiliation-zh = {机电工程学院},
  major-zh      = {机械制造及其自动化},
  student-id    = {□□□□□□□□□□},
} }
\ExplSyntaxOn
%    \end{macrocode}
%
%    \begin{macrocode}
%</hithesis-cfg>
%    \end{macrocode}
%
% \Finale
